% FORMALIZACION_NUEVA: prueba explícita de la descomposición y la
% irreducibilidad ya enunciadas en el desarrollo pentacomponente de REV06.
% CERTIFICADO_NUEVO: verificar_caracteres_pentadicos_rev07.py.
% Inserción prevista tras la definición de V12, antes de su descomposición.
\paragraph{Caracteres de la incidencia icosaédrica.}
El cálculo siguiente se efectúa sobre la acción icosaédrica obtenida en el
teorema~\ref{ii:pent:thm:gravity-exceptional-descent}, con sus hipótesis
de incidencia y orientación. Determina la descomposición de la
representación de vértices y la irreducibilidad del espacio tensorial
mediante sus trazas, conservando el antecedente
\(\pi_{\mathrm{ico}}\) declarado en ese teorema.

\begin{proposition}[Descomposición de vértices y carácter tensorial]
\label{prop:ii-pent-caracteres}
Sean \(\tau=(1+\sqrt5)/2\) y
\(\bar\tau=(1-\sqrt5)/2\). En el orden de clases
\(1,2A,3A,5A,5B\), los caracteres de la representación rotacional
\(V_3\), de su conjugada algebraica \(V_{3'}\), de la representación
de los doce vértices \(V_{12}\) y de
\(\operatorname{STF}_3=\operatorname{Sym}^2_0(V_3)\) son
\begin{equation}
\label{eq:ii-pent-tabla-caracteres}
\begin{array}{c|rrrrr}
\text{clase} &1&2A&3A&5A&5B\\ \hline
\text{cardinal}&1&15&20&12&12\\
\text{clase de }g^2&1&1&3A&5B&5A\\ \hline
\chi_3&3&-1&0&\tau&\bar\tau\\
\chi_{3'}&3&-1&0&\bar\tau&\tau\\
\chi_{12}&12&0&0&2&2\\
\chi_{\mathrm{STF}}&5&1&-1&0&0
\end{array}
\end{equation}
Las representaciones \(V_3,V_{3'}\) y
\(\operatorname{STF}_3\) son absolutamente irreducibles y distintas.
Además,
\begin{equation}
\label{eq:ii-pent-descomposicion-demostrada}
 V_{12}\simeq\mathbf1\oplus V_3\oplus V_{3'}
                  \oplus\operatorname{STF}_3,
 \qquad
 \langle\chi_{12},\chi_{\mathrm{STF}}\rangle=1.
\end{equation}
En particular, el sumando de dimensión cinco aparece una sola vez.
\end{proposition}

\begin{proof}
En \(A_5\), las dobles transposiciones son \(5\cdot3=15\) y
los ciclos de longitud tres son \(\binom53\cdot2=20\).
Los \(4!=24\) ciclos de longitud cinco se separan en dos clases de
doce: su centralizador es el subgrupo cíclico de orden cinco y está
contenido en \(A_5\). Si se identifican las posiciones de un ciclo
con \(\mathbb Z/5\mathbb Z\), la permutación que induce
\(r\mapsto2r\) es un ciclo de longitud cuatro sobre los elementos
no nulos, de signo negativo. Por ello el cuadrado intercambia ambas
clases. La inversión \(r\mapsto-r\) es producto de dos
transposiciones y conserva cada clase. Esto determina las dos primeras
filas de~\eqref{eq:ii-pent-tabla-caracteres}.

La representación \(V_3\) es la acción por rotaciones sobre las
coordenadas de los vértices. Una rotación de ángulo \(\theta\)
tiene traza \(1+2\cos\theta\). Las clases de órdenes dos y tres
dan, respectivamente, \(-1\) y \(0\). Se denomina \(5A\) a la
clase de ángulos \(\pm2\pi/5\) y \(5B\) a la de
\(\pm4\pi/5\); sus trazas son \(\tau\) y \(\bar\tau\).
Estas expresiones se obtienen de
\(u^2+u-1=0\), para \(u=2\cos(2\pi/5)\), y de la fórmula
del ángulo doble. La ecuación de \(u\) resulta al dividir
\(1+z+z^2+z^3+z^4=0\) por \(z^2\), con
\(z=\exp(2\pi i/5)\). Se ha calculado así el carácter de una
representación ya definida geométricamente.

Para construir la otra representación, se usan los vértices sin la
normalización común \(\nu^{-1}\). Sus coordenadas pertenecen a
\(K=\mathbb Q(\sqrt5)\). Si \(B\) es la matriz de tres vértices
linealmente independientes y \(B_g\) contiene sus imágenes, la
matriz rotacional es \(R_g=B_gB^{-1}\in M_3(K)\).
La conjugación \(\sigma:\sqrt5\mapsto-\sqrt5\) define entonces
\(R'_g=\sigma(R_g)\), y satisface
\[
 R'_{gh}=R'_gR'_h,\qquad
 (R'_g)^TR'_g=I,\qquad \det R'_g=1.
\]
Por tanto, \(V_{3'}\) es una representación real efectivamente
construida. Su carácter es \(\sigma(\chi_3)\), que intercambia
\(\tau\) y \(\bar\tau\).

El carácter de permutación cuenta vértices fijos. La identidad fija
doce; las rotaciones no triviales de orden cinco fijan los dos extremos
de su eje. Las de orden dos o tres no fijan vértices, pues el
estabilizador de un vértice es \(C_5\). Se obtiene así
\(\chi_{12}=(12,0,0,2,2)\).

Para el espacio tensorial, si \(\lambda_1,\lambda_2,\lambda_3\)
son los autovalores de \(R_g\), la suma
\(\sum_{i\leq j}\lambda_i\lambda_j\) es
\(\bigl((\sum_i\lambda_i)^2+\sum_i\lambda_i^2\bigr)/2\).
La forma métrica genera una recta invariante en
\(\operatorname{Sym}^2(V_3)\), cuyo complemento sin traza tiene
por carácter
\begin{equation}
\label{eq:ii-pent-caracter-stf}
 \chi_{\mathrm{STF}}(g)
 =\frac{\chi_3(g)^2+\chi_3(g^2)}2-1.
\end{equation}
Al usar el mapa de cuadrados de la tabla y
\(\tau+\bar\tau=1\), \(\tau\bar\tau=-1\), resulta
\(\chi_{\mathrm{STF}}=(5,1,-1,0,0)\).

El producto de caracteres reales es
\(\langle\chi,\psi\rangle
=\frac1{60}\sum_C|C|\chi(C)\psi(C)\).
En particular,
\begin{align*}
 \langle\chi_3,\chi_3\rangle
 &=\langle\chi_{3'},\chi_{3'}\rangle
 =\frac{9+15+12(\tau^2+\bar\tau^2)}{60}=1,\\
 \langle\chi_3,\chi_{3'}\rangle
 &=\frac{9+15+24\tau\bar\tau}{60}=0,\\
 \langle\chi_{\mathrm{STF}},\chi_{\mathrm{STF}}\rangle
 &=\frac{25+15+20}{60}=1.
\end{align*}
La norma uno prueba irreducibilidad de las complejificaciones; por
consiguiente, también irreducibilidad absoluta de las representaciones
reales construidas. Sus dimensiones y la ortogonalidad distinguen los
tres sumandos. Finalmente, la igualdad de funciones de clase
\[
 \chi_{12}=1+\chi_3+\chi_{3'}+\chi_{\mathrm{STF}}
\]
se verifica en las cinco columnas. La semisimplicidad de las
representaciones de un grupo finito en característica cero da la
descomposición~\eqref{eq:ii-pent-descomposicion-demostrada}. La
multiplicidad del sumando tensorial puede comprobarse directamente:
\[
 \langle\chi_{12},\chi_{\mathrm{STF}}\rangle
 =\frac{12\cdot5}{60}=1.
\]
\end{proof}

Así, el idempotente central asociado a
\(\chi_{\mathrm{STF}}=\chi_5\) tiene rango
\(5\cdot1=5\) en \(V_{12}\). La irreducibilidad que interviene
en la normalización del entrelazador tensorial y la multiplicidad que
interviene en el acoplamiento gravitatorio quedan acreditadas por
este mismo cálculo.
